%------------------------
% Curriculum Vitae — Junseok Oh (English)
% Compile: xelatex cv_en.tex
%------------------------

\documentclass[a4paper,9pt]{extarticle}

\usepackage{kotex}
\usepackage{fontspec}

\newenvironment{kotext}{%
  \begingroup
  \setmainfont{Pretendard}[
    BoldFont = {Pretendard Bold},
    ItalicFont = {Pretendard Regular},
    BoldItalicFont = {Pretendard Bold},
  ]
  \setmainhangulfont{Pretendard}[
    BoldFont = {Pretendard Bold},
    ItalicFont = {Pretendard Regular},
    BoldItalicFont = {Pretendard Bold},
  ]
}{%
  \endgroup
}
\newcommand{\koline}[1]{\begin{kotext}#1\end{kotext}}
\usepackage{geometry}
\geometry{a4paper, top=0.85in, bottom=0.85in, left=0.9in, right=0.9in}
\usepackage{tabularx}
\usepackage{array}
\usepackage{titlesec}
\usepackage{enumitem}
\usepackage{hyperref}
\usepackage{microtype}

\hypersetup{hidelinks}

\newlength{\cvdatewidth}
\setlength{\cvdatewidth}{3.8cm}

\newenvironment{tightitems}{%
  \begin{itemize}[leftmargin=1.4em, itemsep=1pt, topsep=2pt, parsep=0pt, partopsep=0pt]
}{%
  \end{itemize}%
  \vspace{0.45em}
}

\newenvironment{blockitems}{%
  \begin{itemize}[leftmargin=1.4em, itemsep=5pt, topsep=4pt, parsep=0pt, partopsep=0pt]
}{%
  \end{itemize}
}

\titleformat{\section}{\LARGE\bfseries\scshape}{}{0em}{}[\vspace{0.15em}\titlerule]
\titlespacing*{\section}{0pt}{1.1\baselineskip}{0.6\baselineskip}
\titleformat{\subsection}{\Large\bfseries\itshape}{}{0em}{}
\titlespacing*{\subsection}{0pt}{0.65\baselineskip}{0.35\baselineskip}

\newcommand{\entrytitle}[1]{{\fontsize{10.5}{13}\selectfont\bfseries #1}}
\newcommand{\me}[1]{\textbf{\underline{#1}}}

\newcommand{\cvheader}[1]{%
  \begin{center}
    {\LARGE\bfseries Junseok Oh (오준석)}\\[0.55em]
    {\normalsize\itshape #1}\\[0.75em]
    (+82) 010-6640-0875 \;$|$\;
    \href{mailto:ohjs@sogang.ac.kr}{ohjs@sogang.ac.kr} \;$|$\;
    \href{https://june-oh.github.io/june-oh/}{june-oh.github.io/june-oh}
  \end{center}
  \vspace{0.5em}
}

\newcommand{\eduentry}[4]{%
  \vspace{1.0em}
  \noindent\entrytitle{#1}\hfill #4\par
  \ifx&#3&\else\noindent\hfill\textit{#3}\par\fi
  \noindent #2\par
  \vspace{0.75em}
}

\newcommand{\researchentry}[4]{%
  \vspace{0.55em}
  \noindent
  \begin{tabularx}{\textwidth}{@{}>{\raggedright\arraybackslash}X >{\raggedleft\arraybackslash}p{\cvdatewidth}@{}}
    \entrytitle{#1} & \textit{#2} \\
  \end{tabularx}
  \ifx&#3&\else\noindent\textit{#3}\par\fi
  \vspace{0.1em}
  \begin{itemize}[leftmargin=1.4em, itemsep=1pt, topsep=2pt, parsep=0pt]
    #4
  \end{itemize}
  \vspace{0.25em}
}

\newcommand{\teachingentry}[5]{%
  \vspace{0.5em}
  \noindent
  \begin{tabularx}{\textwidth}{@{}>{\raggedright\arraybackslash}X >{\raggedleft\arraybackslash}p{\cvdatewidth}@{}}
    \entrytitle{#1} & \textit{#2} \\
  \end{tabularx}
  \noindent\textit{#3} \textbar{} #4\par
  \noindent #5\par
  \vspace{0.35em}
}

\begin{document}
\pagestyle{empty}

\cvheader{Speech AI Engineer \textbar{} End-to-End ASR \& Speech-LLM}

\section*{Education}
\eduentry{Sogang University}{Ph.D.\ in Computer Science and Engineering (Advisor: Prof.\ Ji-Hwan Kim)}{Seoul, South Korea}{Mar 2022 -- Aug 2026}
\begin{tightitems}
  \item Dissertation: ``Adapter-Only Bridging of Frozen Speech Encoder and Frozen LLM for ASR''
\end{tightitems}
\eduentry{Sogang University}{Master of Engineering in Computer Science and Engineering (Advisor: Prof.\ Ji-Hwan Kim)}{Seoul, South Korea}{Sep 2017 -- Aug 2019}
\begin{tightitems}
  \item Thesis: ``Korean Real-Time Automatic Transcription System Using Weakly Labeled Corpus''
\end{tightitems}
\eduentry{Sogang University}{Bachelor of Engineering in Computer Science and Engineering}{Seoul, South Korea}{Mar 2010 -- Aug 2017}
\eduentry{Uijeongbu High School}{High School Diploma}{Uijeongbu, Gyeonggi-do, South Korea}{Mar 2007 -- Feb 2010}

\section*{Research Experience}
\researchentry{Development of Dialog-based Multi-modal Explainable AI Technology}{Apr 2022 -- Present}{Funding Agency: Ministry of Science and ICT (MSIT) / IITP}{
  \item Developed an AI-based framework for dysarthria (speech disorder) severity classification, providing multi-modal explanations to support diagnostic decision-making.
  \item Implemented speech-based explainable diagnostic modules that analyze acoustic and linguistic features to provide interpretable feedback on severity levels.
  \item Researched methods to bridge the gap between complex AI model outputs and user-understandable explanations in a multi-modal dialog environment.
}

\researchentry{End-to-End Korean Speech Recognition Project}{2024 -- Apr 2025}{Industry Partner: Smilegate}{
  \item Developed a universal Korean ASR system using a hybrid FastConformer RNN-Transducer (RNNT) + CTC model.
  \item Implemented a cache-aware streaming method to improve recognition accuracy and reduce latency.
  \item Implemented context biasing to adapt the ASR decoder for domain-specific vocabulary, such as gaming terminology, thereby improving recognition accuracy and robustness in specialized contexts.
}

\researchentry{End-to-End Speech Recognition for Telephony}{Apr 2024 -- Dec 2024}{Industry Partner: LOTTE INNOVATE COMPANY}{
  \item Developed streaming and non-streaming Korean ASR pipelines optimized for 8\,kHz telephony data using a FastConformer-CTC architecture.
  \item Implemented context-biasing modules to dynamically adapt the decoder's vocabulary for domain shift, boosting word-level accuracy.
}

\researchentry{Automated Korean Speaking Assessment (2024)}{May 2024 -- Dec 2024}{Funding Agency: Ministry of Culture, Sports and Tourism of South Korea}{
  \item Developed a multi-task learning framework using Wav2Vec to jointly model pronunciation, fluency, and content for L2-Korean assessment.
  \item Integrated Conformer-CTC ASR outputs with a large language model (LLM) to generate automated, multi-aspect scoring.
}

\researchentry{Intelligent Analysis and Classification-based Content Rating Technology to Address Uncontrolled Distribution of Harmful Media}{2022 -- 2024}{Funding Agency: Ministry of Science and ICT of South Korea}{
  \item Led the audio analytics submodule within an automated video content rating framework, designing sound-event detection models for key multimedia events.
  \item Fine-tuned OpenAI's Whisper ASR model on domain-specific video corpora to enhance transcription robustness in diverse acoustic environments.
}

\researchentry{Automated Korean Speaking Assessment (2023)}{May 2023 -- Dec 2023}{Funding Agency: Ministry of Culture, Sports and Tourism of South Korea}{
  \item Built an end-to-end evaluation pipeline for L2 Korean speakers by combining Conformer-CTC ASR outputs with BERT-based semantic scoring.
  \item Developed algorithms to quantify pronunciation accuracy, speech rate, and syntactic correctness from ASR transcripts.
}

\researchentry{Development of a Video Story Understanding-based QA System to Pass the Video Turing Test}{Sep 2017 -- Dec 2019}{Funding Agency: Ministry of Science and ICT of South Korea}{
  \item Modified Kaldi's sentence-level decoder to achieve sub-1.0 RT for real-time video QA applications as part of a four-member research team.
  \item Collected and curated domain-specific audio/text corpora from target video series to optimize acoustic and language models.
  \item Conducted acoustic model retraining and language model updates using the new corpora, improving QA accuracy on complex video narratives.
}

\section*{Publications}
\subsection*{International Journal}
\begin{blockitems}
  \item \me{J. Oh}, J. Nam, and J.-H. Kim, ``HiTCA: Fusing Hierarchical Text and Contextual Audio for Accurate VCR,'' \textit{EURASIP Journal on Audio, Speech, and Music Processing}, under review. [SCIE]
  \item S. Ma, \me{J. Oh}, M. Kim, and J.-H. Kim, ``Survey on Deep Learning-based Speech Technologies in Voice Chatbot Systems,'' \textit{KSII Transactions on Internet \& Information Systems (TIIS)}, vol.~19, no.~5, pp.~1406--1440, 2025. [SCIE]
  \item \me{J. Oh}, E. Cho, and J.-H. Kim, ``Integration of WFST language model in pre-trained Korean E2E ASR model,'' \textit{KSII Transactions on Internet and Information Systems (TIIS)}, vol.~18, no.~6, pp.~1692--1705, 2024. [SCIE]
  \item S. Seo, \me{J. Oh}, E. Cho, H. Park, G. Kim, and J.-H. Kim, ``TP-MobNet: A Two-pass Mobile Network for Low-complexity Classification of Acoustic Scene,'' \textit{Computers, Materials \& Continua}, vol.~73, no.~2, 2022. [SCIE]
  \item M. Lim, D. Lee, H. Park, Y. Kang, \me{J. Oh}, J.-S. Park, G.-J. Jang, and J.-H. Kim, ``Convolutional neural network based audio event classification,'' \textit{KSII Transactions on Internet and Information Systems (TIIS)}, vol.~12, no.~6, pp.~2748--2760, Jun. 2018. [SCIE]
\end{blockitems}

\subsection*{International Conference}
\begin{blockitems}
  \item S. Ma, \me{J. Oh}, and J.-H. Kim, ``Retention-Preserving Gradient Projection with Entropy-Guided Token-Level Distillation for Rehearsal-Free Continual ASR,'' in \textit{Proc.\ Interspeech}, Sydney, Australia, 2026. [Accepted]
  \item \me{J. Oh} and J.-H. Kim, ``Adapter-Only Bridging of Frozen Speech Encoder and Frozen LLM for ASR,'' in \textit{Advances in Knowledge Discovery and Data Mining (PAKDD 2026)}, Lecture Notes in Computer Science, vol.~16600, pp.~251--262, Springer, Singapore, 2026. [Oral]
  \item \me{J. Oh} and J.-H. Kim, ``SEAM: Bridging the Temporal-Semantic Granularity Gap for LLM-based Speech Recognition,'' in \textit{Findings of the Association for Computational Linguistics: EACL 2026}, pp.~2135--2144, 2026.
  \item \me{J. Oh}, H. Park, and J.-H. Kim, ``Speech Intelligibility Prediction of Dysarthria Using Deep Convolutional Networks,'' in \textit{Proc.\ APIC-IST 2023}, pp.~236--237, 2023.
  \item M. Kim, \me{J. Oh}, and J.-H. Kim, ``Automated Dysarthria Severity Classification Using Diadochokinetic test and Speech Intelligibility Based on LightGBM,'' in \textit{Proc.\ APIC-IST 2023}, pp.~12--13, 2023.
  \item S. Seo, M. Lim, D. Lee, H. Park, \me{J. Oh}, D. J. Rim, and J.-H. Kim, ``Environmental noise robustness for Korean fricatives using speech enhancement generative adversarial networks,'' in \textit{Proc.\ IEEE Int.\ Conf.\ Big Data and Smart Computing (BigComp)}, pp.~1--4, 2019.
  \item S. Seo, D. J. Rim, M. Lim, D. Lee, H. Park, \me{J. Oh}, C. Kim, and J.-H. Kim, ``Shortcut connections based deep speaker embeddings for end-to-end speaker verification system,'' in \textit{Proc.\ Interspeech}, pp.~2928--2932, 2019.
\end{blockitems}

\subsection*{Domestic Journal}
\begin{blockitems}
  \item \koline{이정필, 장재후, 김지현, 김민섭, 김성준, 김민서, 김하영, \me{오준석}, 정원, 김장연 외, ``음성에 기반한 마비말장애 진단과 설명이 가능한 시스템,'' \textit{정보과학회지}, 42(4): 45--56, 2024. [KSCI/KCI]}
  \item H. Park, Y. Kang, M. Lim, D. Lee, \me{J. Oh}, and J.-H. Kim, ``LFMMI-based acoustic modeling by using external knowledge,'' \textit{The Journal of the Acoustical Society of Korea}, vol.~38, no.~5, pp.~607--613, 2019. [KSCI/KCI]
\end{blockitems}

\section*{Awards}
\begin{blockitems}
  \item \koline{\textbf{2023 한국어 AI 경진대회} (NIA) --- Track2-1, 상담 음성인식: 오준석, 김민서, 남주형 (team ``상담ONE''), 장려상}
  \item \koline{\textbf{2022 한국어 인공지능 경진대회} (NIA) --- 기업현안(회의음성): 오준석, 김하영 (team ``SGCSE''), 최우수상/네이버 대표 (1위)}
  \item \koline{\textbf{2021 음절인식률 측정 알고리즘 개발 대회} (KT alpha) --- 박호성, 오준석, 조은수 (team ``검은사케동''), 최우수상 (1위)}
\end{blockitems}

\section*{Teaching}
\teachingentry{Audio Recognition, Synthesis \& Transformation}{Fall 2024}{Teaching Assistant}{CSE5109/CSEG109/AIEG109/AIE5109 \textbar{} Sogang University \textbar{} Prof.\ Ji-Hwan Kim}{Led lab sessions on audio processing, deep learning basics, language models, and FastSpeech2 TTS using PyTorch and Colab.}
\teachingentry{Deep Learning-based Automatic Speech Recognition}{Summer 2023}{Teaching Assistant}{Samsung AI Academy \textbar{} Sogang University $\times$ Samsung Electronics \textbar{} Prof.\ Ji-Hwan Kim}{Hands-on ASR tutorial covering audio handling, MLP, CTC, Whisper, NVIDIA NeMo fine-tuning, and WFST using PyTorch and Colab notebooks.}
\teachingentry{Audio Recognition, Synthesis \& Transformation}{Fall 2023}{Teaching Assistant}{CSE5109/CSEG109/AIEG109/AIE5109 \textbar{} Sogang University \textbar{} Prof.\ Ji-Hwan Kim}{Led lab sessions on audio processing, PyTorch, RNN/CNN/Seq2Seq, and FastSpeech2/VocGAN TTS using PyTorch and Colab.}
\teachingentry{Conversational User Interface}{Spring/Fall 2022}{Teaching Assistant}{CSE5311/CSEG311/GITA370 \textbar{} Sogang University \textbar{} Prof.\ Ji-Hwan Kim}{Led lab sessions on dialogue systems and conversational AI interface design.}

\section*{Patents}
\begin{blockitems}
  \item KR 10-2699607 (B1) --- \textit{Corpus Construction Service Provision Server and Method} \koline{(코퍼스 구축 서비스 제공 서버 및 방법).}\\
  Inventors: Ji-Hwan Kim; Junseok Oh; Hosung Park.\\
  Assignee: Sogang University Industry--Academic Cooperation Foundation.\\
  Filed: Apr 29, 2020 (KR 10-2020-0052570); Published: Nov 8, 2021 (KR 10-2021-0133667); Granted: Aug 22, 2024; Gazette: Aug 26, 2024.
\end{blockitems}

\section*{Certificates}
\begin{blockitems}
  \item NVIDIA Deep Learning Institute Certificate, \textit{Building Conversational AI Applications}, 2022
\end{blockitems}

\section*{Technical Skills}
\noindent
\begin{tabular}{@{}p{0.46\textwidth}@{\hspace{0.04\textwidth}}p{0.46\textwidth}@{}}
  \textbf{Programming Languages \& Scripting} & Python \\
  \textbf{Deep Learning Frameworks} & PyTorch, NVIDIA NeMo, Hugging Face Transformers, PEFT \\
  \textbf{E2E ASR Toolkits} & Kaldi \\
  \textbf{Other Tools \& Libraries} & KenLM \\
\end{tabular}

\end{document}
